% Options for packages loaded elsewhere
\PassOptionsToPackage{unicode}{hyperref}
\PassOptionsToPackage{hyphens}{url}
\documentclass[
]{article}
\usepackage{xcolor}
\usepackage[margin=1in]{geometry}
\usepackage{amsmath,amssymb}
\setcounter{secnumdepth}{-\maxdimen} % remove section numbering
\usepackage{iftex}
\ifPDFTeX
  \usepackage[T1]{fontenc}
  \usepackage[utf8]{inputenc}
  \usepackage{textcomp} % provide euro and other symbols
\else % if luatex or xetex
  \usepackage{unicode-math} % this also loads fontspec
  \defaultfontfeatures{Scale=MatchLowercase}
  \defaultfontfeatures[\rmfamily]{Ligatures=TeX,Scale=1}
\fi
\usepackage{lmodern}
\ifPDFTeX\else
  % xetex/luatex font selection
\fi
% Use upquote if available, for straight quotes in verbatim environments
\IfFileExists{upquote.sty}{\usepackage{upquote}}{}
\IfFileExists{microtype.sty}{% use microtype if available
  \usepackage[]{microtype}
  \UseMicrotypeSet[protrusion]{basicmath} % disable protrusion for tt fonts
}{}
\makeatletter
\@ifundefined{KOMAClassName}{% if non-KOMA class
  \IfFileExists{parskip.sty}{%
    \usepackage{parskip}
  }{% else
    \setlength{\parindent}{0pt}
    \setlength{\parskip}{6pt plus 2pt minus 1pt}}
}{% if KOMA class
  \KOMAoptions{parskip=half}}
\makeatother
\usepackage{color}
\usepackage{fancyvrb}
\newcommand{\VerbBar}{|}
\newcommand{\VERB}{\Verb[commandchars=\\\{\}]}
\DefineVerbatimEnvironment{Highlighting}{Verbatim}{commandchars=\\\{\}}
% Add ',fontsize=\small' for more characters per line
\usepackage{framed}
\definecolor{shadecolor}{RGB}{248,248,248}
\newenvironment{Shaded}{\begin{snugshade}}{\end{snugshade}}
\newcommand{\AlertTok}[1]{\textcolor[rgb]{0.94,0.16,0.16}{#1}}
\newcommand{\AnnotationTok}[1]{\textcolor[rgb]{0.56,0.35,0.01}{\textbf{\textit{#1}}}}
\newcommand{\AttributeTok}[1]{\textcolor[rgb]{0.13,0.29,0.53}{#1}}
\newcommand{\BaseNTok}[1]{\textcolor[rgb]{0.00,0.00,0.81}{#1}}
\newcommand{\BuiltInTok}[1]{#1}
\newcommand{\CharTok}[1]{\textcolor[rgb]{0.31,0.60,0.02}{#1}}
\newcommand{\CommentTok}[1]{\textcolor[rgb]{0.56,0.35,0.01}{\textit{#1}}}
\newcommand{\CommentVarTok}[1]{\textcolor[rgb]{0.56,0.35,0.01}{\textbf{\textit{#1}}}}
\newcommand{\ConstantTok}[1]{\textcolor[rgb]{0.56,0.35,0.01}{#1}}
\newcommand{\ControlFlowTok}[1]{\textcolor[rgb]{0.13,0.29,0.53}{\textbf{#1}}}
\newcommand{\DataTypeTok}[1]{\textcolor[rgb]{0.13,0.29,0.53}{#1}}
\newcommand{\DecValTok}[1]{\textcolor[rgb]{0.00,0.00,0.81}{#1}}
\newcommand{\DocumentationTok}[1]{\textcolor[rgb]{0.56,0.35,0.01}{\textbf{\textit{#1}}}}
\newcommand{\ErrorTok}[1]{\textcolor[rgb]{0.64,0.00,0.00}{\textbf{#1}}}
\newcommand{\ExtensionTok}[1]{#1}
\newcommand{\FloatTok}[1]{\textcolor[rgb]{0.00,0.00,0.81}{#1}}
\newcommand{\FunctionTok}[1]{\textcolor[rgb]{0.13,0.29,0.53}{\textbf{#1}}}
\newcommand{\ImportTok}[1]{#1}
\newcommand{\InformationTok}[1]{\textcolor[rgb]{0.56,0.35,0.01}{\textbf{\textit{#1}}}}
\newcommand{\KeywordTok}[1]{\textcolor[rgb]{0.13,0.29,0.53}{\textbf{#1}}}
\newcommand{\NormalTok}[1]{#1}
\newcommand{\OperatorTok}[1]{\textcolor[rgb]{0.81,0.36,0.00}{\textbf{#1}}}
\newcommand{\OtherTok}[1]{\textcolor[rgb]{0.56,0.35,0.01}{#1}}
\newcommand{\PreprocessorTok}[1]{\textcolor[rgb]{0.56,0.35,0.01}{\textit{#1}}}
\newcommand{\RegionMarkerTok}[1]{#1}
\newcommand{\SpecialCharTok}[1]{\textcolor[rgb]{0.81,0.36,0.00}{\textbf{#1}}}
\newcommand{\SpecialStringTok}[1]{\textcolor[rgb]{0.31,0.60,0.02}{#1}}
\newcommand{\StringTok}[1]{\textcolor[rgb]{0.31,0.60,0.02}{#1}}
\newcommand{\VariableTok}[1]{\textcolor[rgb]{0.00,0.00,0.00}{#1}}
\newcommand{\VerbatimStringTok}[1]{\textcolor[rgb]{0.31,0.60,0.02}{#1}}
\newcommand{\WarningTok}[1]{\textcolor[rgb]{0.56,0.35,0.01}{\textbf{\textit{#1}}}}
\usepackage{graphicx}
\makeatletter
\newsavebox\pandoc@box
\newcommand*\pandocbounded[1]{% scales image to fit in text height/width
  \sbox\pandoc@box{#1}%
  \Gscale@div\@tempa{\textheight}{\dimexpr\ht\pandoc@box+\dp\pandoc@box\relax}%
  \Gscale@div\@tempb{\linewidth}{\wd\pandoc@box}%
  \ifdim\@tempb\p@<\@tempa\p@\let\@tempa\@tempb\fi% select the smaller of both
  \ifdim\@tempa\p@<\p@\scalebox{\@tempa}{\usebox\pandoc@box}%
  \else\usebox{\pandoc@box}%
  \fi%
}
% Set default figure placement to htbp
\def\fps@figure{htbp}
\makeatother
\setlength{\emergencystretch}{3em} % prevent overfull lines
\providecommand{\tightlist}{%
  \setlength{\itemsep}{0pt}\setlength{\parskip}{0pt}}
\usepackage{bookmark}
\IfFileExists{xurl.sty}{\usepackage{xurl}}{} % add URL line breaks if available
\urlstyle{same}
\hypersetup{
  pdftitle={Adaptive{,} non-inferiority{,} equivalence trials},
  hidelinks,
  pdfcreator={LaTeX via pandoc}}

\title{Adaptive, non-inferiority, equivalence trials}
\author{}
\date{\vspace{-2.5em}}

\begin{document}
\maketitle

\textbf{Inference lab} using the five-step template: Hypothesis →
Visualise → Assumptions → Conduct → Conclude.

\subsection{Learning objectives}\label{learning-objectives}

\begin{itemize}
\tightlist
\item
  Distinguish superiority, non-inferiority, and equivalence trials by
  their null hypotheses.
\item
  Choose a non-inferiority margin and say what it means clinically.
\item
  Run and plot a non-inferiority test with a confidence interval versus
  the margin.
\end{itemize}

\subsection{Prerequisites}\label{prerequisites}

Inferential statistics and confidence intervals.

\subsection{Background}\label{background}

A superiority trial asks whether a new treatment beats the comparator. A
non-inferiority (NI) trial asks whether the new treatment is \emph{not
worse than} the comparator by more than a pre-specified margin (often
called Δ). An equivalence trial asks whether the new treatment is
neither better nor worse than the comparator by more than Δ on either
side. These are distinct statistical questions and they require
different designs, analyses, and sample sizes.

The NI margin is the hardest thing about an NI trial. It has to be small
enough that clinicians and patients would trade it for the benefits of
the new treatment (lower cost, fewer side effects, easier
administration), but not so small that the trial is infeasible.
Regulators will often ask for a margin no larger than a defined fraction
of the historical treatment effect of the comparator versus placebo.

Adaptive trials modify some aspect of the design (sample size, arm
allocation, the decision rule) during the trial based on interim
results, using pre-specified rules. Group sequential designs, for
example, stop early for overwhelming efficacy or for futility, while
spending Type I error according to a chosen alpha-spending function.

A one-sided 97.5\% confidence interval that stays on the correct side of
the margin is the operational test in most NI trials. Graphically, this
is a forest plot where you check that the interval does not cross the
margin line.

\subsection{Setup}\label{setup}

\begin{Shaded}
\begin{Highlighting}[]
\FunctionTok{library}\NormalTok{(tidyverse)}
\FunctionTok{set.seed}\NormalTok{(}\DecValTok{42}\NormalTok{)}
\FunctionTok{theme\_set}\NormalTok{(}\FunctionTok{theme\_minimal}\NormalTok{(}\AttributeTok{base\_size =} \DecValTok{12}\NormalTok{))}
\end{Highlighting}
\end{Shaded}

\subsection{1. Hypothesis}\label{hypothesis}

H0: the new treatment is worse than the comparator by more than Δ = 3
units on a continuous outcome. H1: the new treatment is not worse by
more than Δ.

\subsection{2. Visualise}\label{visualise}

\begin{Shaded}
\begin{Highlighting}[]
\NormalTok{trial }\OtherTok{\textless{}{-}} \FunctionTok{tibble}\NormalTok{(}
  \AttributeTok{arm =} \FunctionTok{rep}\NormalTok{(}\FunctionTok{c}\NormalTok{(}\StringTok{"comparator"}\NormalTok{, }\StringTok{"new"}\NormalTok{), }\AttributeTok{each =}\NormalTok{ n),}
  \AttributeTok{y   =} \FunctionTok{c}\NormalTok{(}\FunctionTok{rnorm}\NormalTok{(n, }\DecValTok{70}\NormalTok{, }\DecValTok{10}\NormalTok{), }\FunctionTok{rnorm}\NormalTok{(n, }\DecValTok{69}\NormalTok{, }\DecValTok{10}\NormalTok{))}
\NormalTok{)}

\NormalTok{trial }\SpecialCharTok{|\textgreater{}}
  \FunctionTok{ggplot}\NormalTok{(}\FunctionTok{aes}\NormalTok{(arm, y, }\AttributeTok{fill =}\NormalTok{ arm)) }\SpecialCharTok{+}
  \FunctionTok{geom\_boxplot}\NormalTok{(}\AttributeTok{alpha =} \FloatTok{0.6}\NormalTok{, }\AttributeTok{colour =} \StringTok{"grey30"}\NormalTok{) }\SpecialCharTok{+}
  \FunctionTok{labs}\NormalTok{(}\AttributeTok{x =} \ConstantTok{NULL}\NormalTok{, }\AttributeTok{y =} \StringTok{"Outcome"}\NormalTok{) }\SpecialCharTok{+}
  \FunctionTok{theme}\NormalTok{(}\AttributeTok{legend.position =} \StringTok{"none"}\NormalTok{)}
\end{Highlighting}
\end{Shaded}

\subsection{3. Assumptions}\label{assumptions}

Independent observations, roughly normal within-arm residuals, and a
pre-specified margin declared before looking at the data. The
\emph{direction} of the test matters: we test whether the lower bound of
the mean difference (new − comparator) is above −Δ.

\subsection{4. Conduct}\label{conduct}

\begin{Shaded}
\begin{Highlighting}[]
\NormalTok{fit}
\end{Highlighting}
\end{Shaded}

\begin{Shaded}
\begin{Highlighting}[]
\NormalTok{ci }\OtherTok{\textless{}{-}}\NormalTok{ fit}\SpecialCharTok{$}\NormalTok{conf.int}
\NormalTok{est }\OtherTok{\textless{}{-}} \FunctionTok{diff}\NormalTok{(}\FunctionTok{rev}\NormalTok{(fit}\SpecialCharTok{$}\NormalTok{estimate))  }\CommentTok{\# new {-} comparator}
\NormalTok{ni\_pass }\OtherTok{\textless{}{-}}\NormalTok{ ci[}\DecValTok{1}\NormalTok{] }\SpecialCharTok{\textgreater{}} \SpecialCharTok{{-}}\NormalTok{delta}
\FunctionTok{tibble}\NormalTok{(}\AttributeTok{estimate =}\NormalTok{ est, }\AttributeTok{low =}\NormalTok{ ci[}\DecValTok{1}\NormalTok{], }\AttributeTok{high =}\NormalTok{ ci[}\DecValTok{2}\NormalTok{], }\AttributeTok{margin =} \SpecialCharTok{{-}}\NormalTok{delta, ni\_pass)}
\end{Highlighting}
\end{Shaded}

\begin{Shaded}
\begin{Highlighting}[]
\FunctionTok{tibble}\NormalTok{(}\AttributeTok{estimate =}\NormalTok{ est, }\AttributeTok{low =}\NormalTok{ ci[}\DecValTok{1}\NormalTok{], }\AttributeTok{high =}\NormalTok{ ci[}\DecValTok{2}\NormalTok{]) }\SpecialCharTok{|\textgreater{}}
  \FunctionTok{ggplot}\NormalTok{(}\FunctionTok{aes}\NormalTok{(}\AttributeTok{x =}\NormalTok{ estimate, }\AttributeTok{y =} \DecValTok{1}\NormalTok{)) }\SpecialCharTok{+}
  \FunctionTok{geom\_point}\NormalTok{(}\AttributeTok{size =} \DecValTok{3}\NormalTok{) }\SpecialCharTok{+}
  \FunctionTok{geom\_errorbarh}\NormalTok{(}\FunctionTok{aes}\NormalTok{(}\AttributeTok{xmin =}\NormalTok{ low, }\AttributeTok{xmax =}\NormalTok{ high), }\AttributeTok{height =} \FloatTok{0.1}\NormalTok{) }\SpecialCharTok{+}
  \FunctionTok{geom\_vline}\NormalTok{(}\AttributeTok{xintercept =} \DecValTok{0}\NormalTok{, }\AttributeTok{linetype =} \StringTok{"dashed"}\NormalTok{, }\AttributeTok{colour =} \StringTok{"grey40"}\NormalTok{) }\SpecialCharTok{+}
  \FunctionTok{geom\_vline}\NormalTok{(}\AttributeTok{xintercept =} \SpecialCharTok{{-}}\NormalTok{delta, }\AttributeTok{colour =} \StringTok{"firebrick"}\NormalTok{) }\SpecialCharTok{+}
  \FunctionTok{labs}\NormalTok{(}\AttributeTok{x =} \StringTok{"Mean difference (new {-} comparator)"}\NormalTok{, }\AttributeTok{y =} \ConstantTok{NULL}\NormalTok{) }\SpecialCharTok{+}
  \FunctionTok{theme}\NormalTok{(}\AttributeTok{axis.text.y =} \FunctionTok{element\_blank}\NormalTok{())}
\end{Highlighting}
\end{Shaded}

\subsection{5. Concluding statement}\label{concluding-statement}

\begin{quote}
The new treatment had a mean difference of \texttt{round(est,\ 2)}
versus comparator (95\% CI: \texttt{round(ci{[}1{]},\ 2)} to
\texttt{round(ci{[}2{]},\ 2)}). With a pre-specified non-inferiority
margin of Δ = 3, the new treatment
\texttt{if\ (ni\_pass)\ "met"\ else\ "did\ not\ meet"} the
non-inferiority criterion because the lower bound
\texttt{if\ (ni\_pass)\ "exceeded"\ else\ "fell\ below"} −Δ.
\end{quote}

\subsubsection{A word on adaptive
designs}\label{a-word-on-adaptive-designs}

Adaptive designs work when the adaptation rules and the alpha spent at
each look are pre-specified. Running an interim analysis with the hope
of extending the trial if it looks promising --- without a pre-
specified rule --- inflates Type I error and destroys the trial's
inferential warranty.

\subsection{Common pitfalls}\label{common-pitfalls}

\begin{itemize}
\tightlist
\item
  Picking the NI margin after seeing the data.
\item
  Using a two-sided superiority test and claiming non-inferiority.
\item
  Running an ``interim look'' without a pre-specified spending rule.
\item
  Treating equivalence and non-inferiority as interchangeable.
\end{itemize}

\subsection{Further reading}\label{further-reading}

\begin{itemize}
\tightlist
\item
  Piaggio G et al.~(2012), \emph{CONSORT extension for non-inferiority
  and equivalence trials}.
\item
  Jennison C, Turnbull BW. \emph{Group Sequential Methods with
  Applications to Clinical Trials}.
\item
  FDA (2016), \emph{Non-Inferiority Clinical Trials to Establish
  Effectiveness}.
\end{itemize}

\subsection{Session info}\label{session-info}

\begin{Shaded}
\begin{Highlighting}[]

\end{Highlighting}
\end{Shaded}

\subsection{See also --- chapter index}\label{see-also-chapter-index}

\begin{itemize}
\tightlist
\item
  \href{../index.Rmd}{Methods in this chapter}
\item
  \href{../../../ch00-find-your-method.Rmd}{Find your method (Chapter
  0)}
\end{itemize}

\end{document}
